\documentclass[12pt,a4paper]{article}
\usepackage{kotex}
\usepackage{geometry}
\usepackage{amsmath,amssymb}
\usepackage{graphicx}
\usepackage{booktabs}
\usepackage{hyperref}
\usepackage{cite}
\usepackage{setspace}
\usepackage{array}
\usepackage{footnote}
\geometry{margin=1in}
\onehalfspacing

\hypersetup{
  colorlinks=true,
  linkcolor=blue,
  citecolor=blue,
  urlcolor=blue
}

\title{
  \textbf{BIMOD-006} \\[0.5em]
  \Large BMNP 경락 가설: \\
  교차배열 단일극성 나노입자 네트워크, \\
  광학 게이팅, 8체질 분류의 통합 이론
}
\author{
  윤덕진 (Solvaram) \\
  \textit{팔체질 연구 플랫폼 (Eight-Constitutional Research Platform)} \\
  \small Claude AI (Anthropic) 협업 이론 구성 \\
  \small \texttt{OSF: https://doi.org/10.17605/OSF.IO/VJF63}
}
\date{Version 1.3 --- 2026년 2월}

\begin{document}
\maketitle

% ─────────────────────────────────────────
\begin{abstract}
BIMOD-006은 \textbf{BMNP 경락 가설}을 제안한다. 생체 자성 나노입자(BMNPs)는 각각
안정적인 \emph{단자구(single magnetic domain, SMD)} 상태를 유지하며 명확한
단일극성을 나타낸다. 이 단자구 입자들은 두 가지 구조를 이룬다:
(1)~도전성 통로로서의 \emph{교차배열 단일극성 사슬}(N--S--N--S\ldots)과
(2)~기능적 터미널로서의 \emph{단극성 뭉침(단극성 절점)}.
이 이중 구조 네트워크가 동아시아 의학 경락 체계의 물리적 기반으로 제안된다.

핵심 증거: (i)~감마프로테오박테리아 SS-5 TEM 이미지에서 검은 입자가
S극 위치에 일관 분포(새로운 BIMOD 관찰); (ii)~각 마그네토솜이 독립적
단자구 나노입자임 확인(Nanoscale Advances, 2020); (iii)~식물 체관부 세포벽에서
BMNP 사슬 검출; (iv)~주입 자성 나노입자가 경락 경로를 따라 이동
\cite{johng2007}; (v)~BIMOD 극성 판별의 광자 의존성 게이팅 확인.

나시온 혈위의 4개 지점에서 표현되는 8체질 표현형은 유전적으로 결정된
8가지 BMNP 네트워크 구조를 반영하는 것으로 제안한다.
이 틀은 검증 가능한 예측을 제공하고 객관적 센서 개발의 기반이 된다.

\textbf{핵심어:} BIMOD, 생체 자성 나노입자, 경락 가설, 단일극성 교차배열,
단자구 나노입자, 팔체질, 나시온 혈위, 광학 게이팅, 마그네토솜, 체질 약리학
\end{abstract}

% ─────────────────────────────────────────
\tableofcontents
\newpage

% ─────────────────────────────────────────
\section{서론}

\subsection{경락의 미해결 문제}
동아시아 전통의학의 경락 체계는 수천 년의 임상 적용에도 불구하고
그 해부학적·물리적 기반이 규명되지 않았다.
결합 조직면, 간질 유체 채널, 신경 경로, 봉한관 체계 등이 후보로 제시되었으나
합의에 이르지 못했다.

본 논문은 새로운 후보를 제안한다: 각각 단자구 단일극성을 가진 BMNPs가
N--S--N--S 교차배열로 조직된 네트워크로, 절점에는 동일 극성 뭉침이 존재한다.
이 구조는 해부학적으로 안정된 경로를 따라 방향성 있는 신호 전달을 위한
물리적 기반을 제공한다.

\subsection{단자구 나노입자와 단일극성 현상}
임계 크기($\sim$25~nm) 이상의 안정적 단자구 상태를 이루는 자성 나노입자는
내부 모든 스핀이 하나의 방향으로 정렬된 단일극성 거동을 나타낸다.
마그네토솜 연구 \cite{kornig2020}는 각 입자가 독립적인 단자구 나노입자이며,
자화 방향 변화는 입자의 물리적 이동이 아닌 내부적 일관 회전으로만
일어남을 확인했다.

BIMOD 판독에서 관찰되는 현상은 이와 일치한다:
영구자석 단면 중앙부($\sim$30\% 범위)에서는 본래 극성이,
최외곽 테두리에서는 얇은 반대극성 띠가 감지되고,
중간 영역은 무극성으로 나타나는 쌍극자 구조와 달리,
나노입자 표본에서는 \emph{전체 면적에 걸쳐 단일 극성이 일관되게 포착}된다.
중간 무극성 영역도, 테두리 반대극성도 존재하지 않는다.
이는 각 BMNP 입자가 쌍극자 분산 없이 단자구 단일극성을 나타낸다는
BIMOD 관찰적 증거다.

\subsection{탐지 방법으로서의 BIMOD 프레임워크}
BIMOD(체외 간섭성 자기-광학 판별법)는 4중 공명 체계로 작동한다:
(1)~국소 영구자석 탐침, (2)~BMNP를 함유한 인체 조직,
(3)~지구 지자기장, (4)~광자 매개 광학 게이팅.
물리적 자석뿐 아니라 자석 사진, 역사적 유화, X선·MRI 이미지로부터도
재현 가능한 극성 판별을 시연했다.
이는 침습적 절차 없이 다양한 생물학적 표본의 BMNP 조직을 조사할 수 있는
최초의 실용적 도구를 제공한다.

\subsection{핵심 관찰: 검은 입자와 S극의 결합}
감마프로테오박테리아 SS-5의 마그네토솜 TEM 공개 이미지
(출처: 위키백과 「자기수용」)에 대한 BIMOD 분석에서,
세 패널(A, B, C) 모두에서 N--S--N--S 엄격한 교번 배열이 확인되었다.
결정적으로, 어둡게 변색된 입자들은 예외 없이 S극 위치를 차지했다.
이 관찰은 입자의 화학-물리적 상태(전자 밀도, 산화도)와
자기 극성 정체성 사이의 결합을 시사한다.

% ─────────────────────────────────────────
\section{이론적 틀: BMNP 경락 가설}

\subsection{제안된 네트워크의 구조적 구성 요소}

\begin{table}[h]
\centering
\caption{네트워크 구성 요소}
\begin{tabular}{>{\bfseries}p{3cm} p{5.5cm} p{5.5cm}}
\toprule
구성 요소 & 구조 & 제안 기능 \\
\midrule
단일극성 교차배열 사슬 &
  N극·S극 단자구 입자가 N--S--N--S 순서로 교번 배열.
  각 입자는 쌍극자가 아닌 단자구 단일극성 상태 &
  방향성 극성 신호 전도;
  입자 이동 없이 제자리 극성 연결로 신호 전달 \\[4pt]
단극성 절점 (S 뭉침) &
  경혈 부위에서 S극 단자구 BMNP의 집적 &
  신호 종결, 증폭, 장기-특이적 변환 \\[4pt]
단극성 절점 (N 뭉침) &
  경혈 부위에서 N극 단자구 BMNP의 집적 &
  신호 수신, 체질 기준선 유지 \\[4pt]
나시온 표면 발현 (4개 지점, 각 N/S) &
  두개 정중선 상의 네 개 양측 대칭 단극성 절점이
  피부 표면에 노출 &
  BIMOD를 통한 8체질 표현형 식별 \\
\bottomrule
\end{tabular}
\end{table}

\subsection{신호 전달 메커니즘}
극성 신호는 나노입자의 물리적 이동에 의해 전달되지 않는다.
조직에 고정된 단자구 입자들이 제자리 극성 직렬 연결을 통해 신호를 전달한다.
두 가지 메커니즘 가능성을 병렬로 고려한다.

\paragraph{가능성 A --- 단일극성 교차 연쇄 신호 전달.}
교번 배열된 단자구 입자들이 이웃 입자와의 자기 상호작용을 통해
연쇄적 극성 상태 변화를 유도한다(마그논 전파의 생물학적 유사체).
각 입자가 단일극자(monopolar)인지 쌍극자인지의 여부는 추가 검토가 필요하다.

\paragraph{가능성 B --- 순차 산화환원 연쇄.}
S극 절점의 선호적 산화 상태(4절 참조)가 전자 전달 연쇄를 구성하며,
이온 채널 조절이 그 출력이 된다.
이 모델은 제자리 극성 직렬 연결에 의한 신호 전달을 지지한다.

두 모델은 상호 배타적이지 않을 수 있다.
어느 메커니즘이 우세한지는 SQUID 기반 실시간 경락 신호 측정으로 검증되어야 한다.

\subsection{팔체질 유전형 가설}
나시온 혈위는 얼굴 정중선을 기준으로 좌우 대칭으로 분포하는
\textbf{8개 위치}로 구성된다(그림~\ref{fig:nasionmap} 참조).
같은 색상으로 표시된 좌우 2개 위치가 하나의 체질축을 이루며,
좌우 극성은 언제나 동일하다.
즉 축~1에서 금양체질(Pulmotonia)은 좌우 양쪽 모두 N극,
목양체질(Hepatonia)은 좌우 양쪽 모두 S극이다.
정중선 자체는 무극성으로 비어 있다.
같은 쪽에 인접한 두 위치는 서로 다른 체질축에 속하므로,
8개 위치로 표현하는 것이 혼동을 방지하는 데 더 직관적이다.

BIMOD에서 확인된 8가지 나시온 극성 표현형은 아래 표와 같이 정리된다.
라틴어 및 한글 체질명과 장부 강세 배열은
8체질의학 창시자 연구소(\url{www.ecmed.org})의 공식 명칭을 따른다.%
\footnote{라틴어 및 한글 체질명(목양체질·금양체질 등)과 장부 강세 배열은
8체질의학 창시자 연구소(\url{www.ecmed.org})의 공식 명칭을 따름.}

\begin{table}[h]
\centering
\caption{8가지 나시온 극성 표현형}
\begin{tabular}{llll}
\toprule
나시온 코드 & 라틴어 공식명 & 한글 공식명 & 주요 강세 장부 \\
\midrule
1S & Hepatonia    & 목양체질 & 간(肝) 최강 \\
1N & Pulmotonia   & 금양체질 & 폐(肺) 최강 \\
2N & Cholecystonia & 목음체질 & 담(膽) 최강 \\
2S & Colonotonia  & 금음체질 & 대장(大腸) 최강 \\
3N & Gastrotonia  & 토양체질 & 위(胃) 최강 \\
3S & Vesicotonia  & 수양체질 & 방광(膀胱) 최강 \\
4N & Pancreotonia & 토음체질 & 췌장(膵臟) 최강 \\
4S & Renotonia    & 수음체질 & 신장(腎臟) 최강 \\
\bottomrule
\end{tabular}
\end{table}

\begin{figure}[h]
\centering
\includegraphics[width=0.62\textwidth]{figures/003nasionmap.jpg}
\caption{나시온 혈위 8개 위치의 극성 배치도.
  얼굴 정중선 좌우로 대칭 분포하며,
  \textbf{같은 색상으로 표시된 좌우 2개 위치가 하나의 체질축}을 구성한다.
  좌우 극성은 언제나 동일하다
  (예: 축~1에서 금양체질은 좌우 모두 N극, 목양체질은 좌우 모두 S극).
  정중선 자체는 무극성으로 비어 있다.
  같은 쪽에 인접한 두 위치는 서로 \emph{다른} 체질축임에 유의할 것.
  파랑: 축~1 (1N/1S); 빨강: 축~3 (3N/3S);
  초록: 축~2 (2N/2S); 검정: 축~4 (4N/4S).}
\label{fig:nasionmap}
\end{figure}

나시온 4개 축은 각각 같은 색상의 좌우 2개 위치로 이루어지며,
같은 축의 두 체질은 동일 위치에서 정반대 극성을 나타낸다:
축~1 (1S$\leftrightarrow$1N), 축~2 (2N$\leftrightarrow$2S),
축~3 (3N$\leftrightarrow$3S), 축~4 (4N$\leftrightarrow$4S).

나시온 혈위 1개 지점의 극성 판별만으로 위치(체질축)와 극성(체질 확정)이
동시에 결정된다.
탐침 자석을 동서 방향으로 45도씩 변화시킬 때 나타나는
체질별 극성 반전 패턴의 정보로도 체질 유형을 추론할 수 있으나,
이는 1차 판별을 보조하는 추가 정보다.%
\footnote{나시온 혈위(Nasion acupoint)와 체질축 구조,
BIMOD 판별 프로토콜은 저자(윤덕진)의 독자적 가설 및 관찰에 근거하며,
\url{www.ecmed.org}의 8체질의학 공식 분류 체계와는 무관함.
본 연구의 이론적 동기는 저자의 블로그
(``8체질과 용신오행의 가설'')에서 시작되었으며,
BIMOD~001--005 논문 시리즈에 상세히 기술되어 있음.}

% ─────────────────────────────────────────
\section{가설을 지지하는 증거}

\subsection{검은 입자--S극 결합 (새로운 관찰)}

\begin{figure}[h]
\centering
\includegraphics[width=0.72\textwidth]{figures/Magnetite_magnetosomes_SS5.png}
\caption{감마프로테오박테리아 SS-5 균주의 자철석 마그네토솜 TEM 이미지
  (패널 A, B, C).
  검게 변색된 고전자밀도 입자는 표~\ref{tab:tem}의 BIMOD 분석에서
  예외 없이 S극 위치에 배정된다.
  \textit{출처:} Lefèvre et al.\ (2012), \textit{ISME J} 6:440--450,
  Wikimedia Commons 경유.
  라이선스: CC~BY~3.0
  (\url{https://en.wikipedia.org/wiki/File:Magnetite_magnetosomes_in_Gammaproteobacteria.png}).}
\label{fig:ss5tem}
\end{figure}

\begin{table}[h]
\centering
\caption{SS-5 마그네토솜 TEM 이미지 BIMOD 극성 판독}
\label{tab:tem}
\begin{tabular}{p{2.5cm} p{5.5cm} p{4.5cm}}
\toprule
패널 & 극성 배열 (검은 입자는 괄호) & 관찰 \\
\midrule
(A) 전체 사슬 &
  SN(S)NSNSN(S)NSN(S)N &
  검은 입자 = S 위치 일치 \\[4pt]
(B) 부분 사슬, 8개 입자 &
  SNSNSN(S)N &
  검은 입자 = S 위치 일치 \\[4pt]
(C) 세로 방향 뷰 &
  N--S (박스: N) &
  N 말단 확인 \\
\bottomrule
\end{tabular}
\end{table}

TEM에서 증가된 전자 밀도는 자철석(Fe$_3$O$_4$)이
마그헤마이트($\gamma$-Fe$_2$O$_3$)로 전환되는 더 높은 산화 상태와 상관된다.
이로부터 다음 가설이 도출된다:

\begin{quote}
\textit{S극 단자구 BMNP 절점은 N극 절점에 비해 선호적으로 산화되거나
전자 밀도가 높은 상태로 유지되며, TEM에서 시각적으로 구별 가능한
대비를 생성한다.}
\end{quote}

추가 표본에서 어두운 입자가 짝수 순번(S극)을 공유하는 패턴이 반복 확인된다면,
TEM 이미지를 BIMOD 없이 BMNP 극성 지도화 도구로 활용할 수 있다.

\subsection{단자구 상태 확인: 기존 과학과의 접점}
\cite{kornig2020}은 마그네토솜 사슬 내 각 입자가 독립적인 단자구 나노입자이며,
사슬 축을 따라 단축 자기 이방성을 보이고,
자화 방향 변화는 입자의 물리적 이동이 아닌 내부적 일관 회전으로만
일어남을 확인했다.
이는 나노입자 전체 면적에서 단일극성이 포착되고
쌍극자 특유의 테두리 반대극성이나 중간 무극성 영역이 존재하지 않는다는
BIMOD 관찰과 구조적으로 일치한다.

\subsection{Johng et al. (2007): 나노입자의 경락 경로 이동}
형광 자성 나노입자를 경혈 LR9에 주입했을 때 같은 경락의 원위 경혈 LR3에서
검출되었고, 비경혈 부위 주입에서는 이런 원위 이동이 나타나지 않았다
\cite{johng2007}.
단, 주입된 외래 나노입자는 주입 부위의 고유 극성 신호 네트워크를
교란했을 가능성이 있다.
이 실험의 의의는 경락 경로를 따르는 물리적 통로의 존재를 확인했다는 점이다.

\subsection{식물 체관부의 BMNP 사슬과 동식물 체질 대응 가능성}
\cite{gorobets2023}은 담배, 감자, 완두콩의 체관부 사관 세포벽에서
BMNP 사슬을 확인했다.
교차배열 단일극성 네트워크 구조가 신경계 진화 이전부터
생명체에 보존된 장거리 신호 통합 전략임을 시사한다.
식물 체질 표현형의 BIMOD 판별이 가능하다면,
동물--식물 간 체질역학적 극성 대응성---특정 식물과 체질 간의 친화 및 길항---을
물리적 기반 위에서 설명하는 가능성이 열린다.

\subsection{지자기 동서 성분과 45도 각도 의존성}
\cite{wang2019}은 동서 지자기 수평 성분을 독립 변수로 제어한 실험에서
뇌 알파파 억제 반응을 확인했다.
\cite{chae2019}은 청색광 조건에서 피험자가 동쪽을 포함한 특정 방향으로
지향함을 보였다.
이 발견들은 동서 지자기 성분이 생물학적으로 활성화된 독립 신호임을 확인하며,
탐침 자석의 45도 회전 단계에 따른 체질별 극성 반전 패턴이라는
BIMOD 관찰의 물리적 배경을 제공한다.

% ─────────────────────────────────────────
\section{산화--극성 결합 가설}

\subsection{물리적 기반}
자철석(Fe$_3$O$_4$)은 Fe$^{2+}$와 Fe$^{3+}$를 모두 함유하며,
표면 산화 시 마그헤마이트($\gamma$-Fe$_2$O$_3$)로 전환되어 Fe$^{3+}$만 남는다.
TEM에서 산화도가 높은 영역은 더 어둡게 나타난다.

\textbf{제안 가설:} 교차배열 사슬에서 S극 단자구 절점은
N극 절점보다 체계적으로 더 산화된 상태로 유지된다.
가능한 원인:
(1)~사슬을 따라 대사 전자 공여체에 대한 차별적 접근;
(2)~지자기 영향 하의 방향성 전자 흐름으로 산화환원 기울기 형성;
(3)~체질별 철 대사 효소 차이로 유전형별 산화 상태 프로파일 생성.

\subsection{검증 가능한 예측}
\begin{itemize}
  \item 어떤 종의 BMNP 사슬 TEM 이미지에서도 어두운 입자들이
        짝수 순번(S극)을 일관되게 차지할 것이다.
  \item 분리된 BMNP 사슬의 XPS에서 S극 위치의
        Fe$^{3+}$/Fe$^{2+}$ 비율이 더 높게 나타날 것이다.
  \item 뫼스바우어 분광법이 교번 입자 사이의 이성질체 이동 차이를
        감지할 것이다.
  \item BIMOD는 종·조직 출처에 관계없이 TEM 이미지의 어두운 입자를
        S극으로 일관되게 배정할 것이다.
\end{itemize}

% ─────────────────────────────────────────
\section{체질 약리학과의 연결}

각 체질의 고유한 장부 강약 배열은 약물 대사 효소의 차별적 발현을 암시한다.
이에 더하여 중요한 원리적 확장을 제안한다:

\begin{quote}
\textit{같은 체질축의 반대 극성 체질---통계적으로 99\%에 달하는 자연 선택적 짝(배우자)---에서
특정 약물의 유효 농도는 상대방에게 부작용 농도로 작용하는
극단적 반비례 차별성이 존재할 수 있다.}
\end{quote}

예: 목음체질(Cholecystonia 2N) $\leftrightarrow$ 금음체질(Colonotonia 2S).
이 상반성은 백신·치료제 개발 시 최적 맞춤 기반을 제공하는 가능성 신호로
해석될 여지가 있다.

구체적 가설:
\begin{itemize}
  \item CYP2C9(이부프로펜 대사): Hepatonia~1S·Cholecystonia~2N·Vesicotonia~3S에서
        고활성 $\to$ 해당 약물 효과적.
  \item CYP2E1(아세트아미노펜 대사): Pulmotonia~1N·Colonotonia~2S·Gastrotonia~3N·Renotonia~4S에서
        고활성 $\to$ 해당 약물 효과적.
  \item 같은 축 반대 극성 체질 쌍에서 상기 두 약물의 효과--부작용 관계가
        역전될 것으로 예측됨.
\end{itemize}

% ─────────────────────────────────────────
\section{미해결 과제 및 연구 우선순위}

\begin{table}[h]
\centering
\caption{연구 우선순위}
\begin{tabular}{p{4.5cm} p{5cm} l}
\toprule
과제 & 제안 방법 & 우선순위 \\
\midrule
검은 TEM 입자가 짝수(S극) 위치에 일관 분포하는가? &
  추가 마그네토솜 TEM 이미지 BIMOD 분석 & 즉각 \\
각 BMNP가 단자구 단일극성인가, 쌍극자인가? &
  SQUID·원자자기 현미경 단일 입자 측정 & 높음 \\
신호 메커니즘: 단일극성 교차 연쇄 vs. 산화환원 연쇄? &
  SQUID 기반 실시간 경락 신호 측정 & 중간 \\
식물 표본이 BIMOD 체질 판독을 제공하는가? &
  식물 조직 사진에 BIMOD 프로토콜 적용 & 높음 \\
XPS·뫼스바우어로 S극 위치 Fe 산화 상태 확인? &
  재료과학 연구실과의 협업 & 중간 \\
45도 탐색이 맹검 시험에서 체질별 반전을 생성하는가? &
  체질 분류된 피험자 대상 통제 각도 실험 & 높음 \\
같은 축 반대 극성 짝이 자유 선택 커플의 99\%인가? &
  커플 모집단 규모 BIMOD 체질 분류 & 높음 \\
\bottomrule
\end{tabular}
\end{table}

% ─────────────────────────────────────────
\section{결론}
BMNP 경락 가설은 단일극성 단자구 입자들의 N--S--N--S 교차배열 사슬과
단극성 말단 절점을 고전적 경락 체계의 물리적 기반으로 제안함으로써
BIMOD 연구 프로그램을 발전시킨다.

TEM 이미지에서 어두운 입자들이 S극 위치를 일관되게 차지한다는 새로운 관찰은
BMNP 극성 정체성을 위한 형태학적 마커가 될 수 있다.

체질적으로 이해된 인간 개인은 지구 지자기장에 맞게 진화에 의해 조율된
자기-광학 간섭계이다.
짝 선택의 극성 상보성, 약물 반응의 체질별 상반성, 식물--동물 간 극성 대응---이 모든 것이
단일 물리적 원리로 통합될 때, 인류의 맞춤 시대를 향한 과학적 토대가 형성된다.

% ─────────────────────────────────────────
\section*{감사의 말}
저자는 12년간의 경험적 관찰과 무상(無償)·공명(共鳴)의 원칙 아래 수행된
판별 봉사에 감사한다.
협업적 이론 개발은 Claude~AI(Anthropic), DeepSeek~AI, Gemini~AI(Google) 및
기타 AI 시스템과 함께 수행되었다.
팔체질 연구 플랫폼 공동체의 지속적 관찰과 헌신에 감사한다.

% ─────────────────────────────────────────
\begin{thebibliography}{99}

\bibitem{yoon2026series}
윤덕진 (2026).
\textit{BIMOD 001--005: BIMOD 프레임워크 기초 논문 시리즈}.
팔체질 연구 플랫폼.
OSF: \url{https://doi.org/10.17605/OSF.IO/VJF63}

\bibitem{yoon2026_500}
윤덕진 (2026).
\textit{BIMOD-500: 생체 자성 나노입자에 기반한 체질 분류 및 응용의 통합 이론}.
팔체질 연구 플랫폼.

\bibitem{johng2007}
Johng,~H.\,M.\ et al.\ (2007).
Use of magnetic nanoparticles to visualize threadlike structures inside
lymphatic vessels of rats.
\textit{Evid.\ Based Complement.\ Alternat.\ Med.}, 4(1), 77--82.

\bibitem{gorobets2019}
Gorobets,~O.\ et al.\ (2019).
Biogenic magnetic nanoparticles in Atlantic salmon olfactory tissue.
\textit{Scientific Reports}.

\bibitem{gorobets2023}
Gorobets,~O.\ et al.\ (2023).
Biogenic magnetic nanoparticles in plant phloem sieve-tube cell walls.
\textit{Plant and Soil}.

\bibitem{kornig2020}
K\"{o}rnig,~A.\ et al.\ (2020).
Stability and magnetic properties of magnetosome chains.
\textit{Nanoscale Advances}.

\bibitem{wang2019}
Wang,~C.\,X.\ et al.\ (2019).
Transduction of the geomagnetic field as evidenced from alpha-band
activity in the human brain.
\textit{eNeuro}, 6(2).

\bibitem{chae2019}
Chae,~K.\,S.\ et al.\ (2019).
Human magnetic sensing and orientation toward magnetic north and east.
\textit{PLoS ONE}.

\bibitem{wikipedia_magnetoreception}
Wikipedia.
\textit{자기수용 (Magnetoreception)}.
\url{https://en.wikipedia.org/wiki/Magnetoreception}

\bibitem{ecmed}
8체질의학 창시자 연구소.
\textit{체질 분류 및 공식 명칭}.
\url{www.ecmed.org}

\end{thebibliography}

\end{document}
